\clearpage
\subsection{\texorpdfstring{$\Tw\A$ 中的正合三角}{TwA 中的正合三角}}\label{sec:3q}
\eref{3.28}中的扭曲复形 $C$ 带有典范态射
\begin{equation}\label{eq:3.30}
\begin{aligned}
\relax[i]&\in\Hom^0_{H(\Tw\A)}(Y_1,C),&i&=\begin{pmatrix}0\\e_{Y_1}\end{pmatrix},\\
[p]&\in\Hom^1_{H(\Tw\A)}(C,Y_0),&p&=(S(e_{Y_0}),0),
\end{aligned}
\end{equation}
\SourcePage{58}
其中 $e_{Y_k}$ 是 $Y_k$ 的恒等态射的链代表（容易看出，具体选择无关紧要）。在 Yoneda 嵌入下，它们成为\stref{引理}{3.6}中的 $\Ai$ 模同态 $\tilde\iota$ 与 $\tilde\pi$。由此以及正合三角和三角化 $\Ai$ 范畴的定义，立即得到：
\begin{lemma}\label{st:3.27}
$H^0(\Tw\A)$ 中的三角\eref{3.9}正合，当且仅当存在 $H^0(\Tw\A)$ 中的同构 $[b]:Y_2\to\Cone(c_1)$，使得
\[
[\mu^2_{\Tw\A}(p,b)]=[c_3],\qquad [i]=[\mu^2_{\Tw\A}(b,c_2)].
\]
\end{lemma}
\begin{lemma}\label{st:3.28}
$\Tw\A$ 是三角化的 $\Ai$ 范畴。
\end{lemma}
构造\eref{3.28}以加性范畴上的链复形范畴中的经典映射锥为原型，在大多数方面也具有相同的性质。特别地，证明其相应同伦范畴为三角范畴的标准论证\bcite{61}（第 IV 章 §2），可以毫无困难地推广到我们的情形，从而得到：
\begin{proposition}\label{st:3.29}
赋予 $H^0(\Tw\A)$ 平移函子 $T=H^0(S)$，并像\stref{命题}{3.14}那样，把正合三角规定为殊三角。这样，$H^0(\Tw\A)$ 就成为经典意义下的三角范畴。
\end{proposition}
由 $\Ai$ 函子 $\F$ 诱导的函子 $\Tw\F:\Tw\A\to\Tw\B$ 与 $S$ 严格交换。此外，可以作等同 $\Tw\F(\Cone(c))\cong\Cone((\Tw\F)^1(c))$，且这一等同与两边相应的态射\eref{3.30}相容。因此：
\begin{lemma}\label{st:3.30}
$H(\Tw\F):H(\Tw\A)\to H(\Tw\B)$ 将正合三角映为正合三角；因而 $H^0(\Tw\F):H^0(\Tw\A)\to H^0(\Tw\B)$ 是经典三角范畴之间的正合函子。
\end{lemma}

\subsection{一般的正合三角}\label{sec:3r}
现在来推导\secref{3g}至\secref{3j}中未经证明而陈述的一般 $\Ai$ 范畴的正合三角性质。主要工具是：
\begin{lemma}\label{st:3.31}
$H(\A)$ 中的三角正合，当且仅当它在嵌入 $H(\A)\hookrightarrow H(\Tw\A)$ 下的像正合。
\end{lemma}
\begin{proof}
取 $\A$ 中的对象 $Y_k$ 和态射 $c_k$，使它们构成一个三角。\stref{引理}{3.7}给出的正合性充要条件，要求存在满足关系\eref{3.10}的某些其他态射，并要求相应的链复形\eref{3.11}无上同调。对于前一个条件，把三角视为位于 $H(\A)$ 还是 $H(\Tw\A)$ 中没有区别，因为嵌入 $\A\hookrightarrow\Tw\A$ 在链层次上就是全忠实的。
\SourcePage{59}
对于后一个条件，若\eref{3.11}对所有 $X\in\Ob\Tw\A$ 都无上同调，那么对所有 $X\in\Ob\A$ 当然也如此。反过来，给定 $X\in\Ob\Tw\A$，可以利用它的子复形递降滤过 $F^iX$，在
\[
\hom_{\Tw\A}(X,Y_2)[1]\oplus\hom_{\Tw\A}(X,Y_0)[1]\oplus\hom_{\Tw\A}(X,Y_1)
\]
上定义一个递增滤过，使之与\eref{3.11}中的微分相容。该滤过的每一个相伴分次部分，都是对 $\A$ 中某些对象取复形\eref{3.11}、再作平移后所得的有限直和。因此，若这些复形均无上同调，与整个 $X$ 相应的复形也无上同调。
\end{proof}
由此很容易得到\stref{推论}{3.15}至\stref{推论}{3.19}所述的正合三角性质：根据\stref{命题}{3.29}，它们在 $H(\Tw\A)$ 中成立，因此在全子范畴 $H(\A)$ 中也成立。

\begin{proof}[\stref{推论}{3.9}的证明]
设 $H(\A)$ 中有正合三角\eref{3.9}，并有 $\Ai$ 函子 $\F:\A\to\B$。由\stref{引理}{3.31}，该三角在 $H(\Tw\A)$ 中也正合。由\stref{引理}{3.30}，它在 $\Tw\F$ 下的像是 $H(\Tw\B)$ 中的正合三角。而 $\Tw\F$ 在 $\A\subset\Tw\A$ 上的限制就是原来的 $\F$；再用一次\stref{引理}{3.31}，即得结论。
\end{proof}

\begin{proof}[\stref{命题}{3.8}的证明（同时也直接证明\stref{推论}{3.10}）]
\stref{引理}{3.7}的一个平凡特例（取 $h_1=h_2=k=0$）表明，$\mathcal D$ 中由对象 $Z_k$ 与态射 $x_k$ 构成的原型三角正合。于是由\stref{推论}{3.9}，对于任意 $\Ai$ 函子 $\F:\mathcal D\to\A$，对象 $Y_k=\F(Z_k)$ 与态射 $[c_k]=[\F^1(x_k)]$ 都构成一个正合三角。

为证明逆向结论，先假设 $\A$ 严格含幺，从而 $\Tw\A$ 也严格含幺。取零次上闭链 $c\in\hom^0_{\A}(Y_0,Y_1)$。考虑非全 $\Ai$ 子范畴 $\widetilde{\mathcal D}\subset\Tw\A$：其对象为 $\widetilde Z_0=Y_0$、$\widetilde Z_1=Y_1$、$\widetilde Z_2=\Cone(c)$，态射空间由 $e_{Y_0},e_{Y_1},c$ 的所有可能出现形式生成：
\begin{align*}
\hom_{\widetilde{\mathcal D}}(\widetilde Z_0,\widetilde Z_0)&=\K e_{Y_0},\\
\hom_{\widetilde{\mathcal D}}(\widetilde Z_1,\widetilde Z_1)&=\K e_{Y_1},\\
\hom_{\widetilde{\mathcal D}}(\widetilde Z_2,\widetilde Z_2)&=\left\{\begin{pmatrix}\lambda_{00}S(e_{Y_0})&0\\\lambda_{10}S(c)&\lambda_{11}e_{Y_1}\end{pmatrix}\right\},\\
\hom_{\widetilde{\mathcal D}}(\widetilde Z_0,\widetilde Z_1)&=\K c,\\
\hom_{\widetilde{\mathcal D}}(\widetilde Z_1,\widetilde Z_2)&=\K\begin{pmatrix}0\\e_{Y_1}\end{pmatrix},\\
\hom_{\widetilde{\mathcal D}}(\widetilde Z_2,\widetilde Z_0)&=\K(S(e_{Y_0}),0),\\
\hom_{\widetilde{\mathcal D}}(\widetilde Z_1,\widetilde Z_0)&=0,\\
\hom_{\widetilde{\mathcal D}}(\widetilde Z_2,\widetilde Z_1)&=\{(\nu_{10}S(c),\nu_{11}e_{Y_1})\},\\
\hom_{\widetilde{\mathcal D}}(\widetilde Z_0,\widetilde Z_2)&=\left\{\begin{pmatrix}\sigma_{00}e_{Y_0}\\\sigma_{10}c\end{pmatrix}\right\}.
\end{align*}
\SourcePage{60}
其中 $\lambda,\nu,\sigma$ 均在 $\K$ 中取值。对于 $d>2$，复合 $\mu^d_{\widetilde{\mathcal D}}$ 为零，因此 $\widetilde{\mathcal D}$ 是 dg 范畴。现在在态射空间 $\hom_{\widetilde{\mathcal D}}(\widetilde Z_i,\widetilde Z_j)$ 上引入自同态 $T^1$。相对于上面给出的基，它们为
\[
\begin{aligned}
T^1(\lambda_{00},\lambda_{10},\lambda_{11})&=(0,0,-\lambda_{10}),\\
T^1(\nu_{10},\nu_{11})&=(0,-\nu_{10}),\\
T^1(\sigma_{00},\sigma_{10})&=(-\sigma_{10},0),
\end{aligned}
\]
其余情形均为零。于是 $\id+\mu^1_{\widetilde{\mathcal D}}T^1+T^1\mu^1_{\widetilde{\mathcal D}}$ 是到一个代表 $\mu^1_{\widetilde{\mathcal D}}$ 上同调的子空间的投影。按照\stref{笔记}{1.13}所述的方法，应用扰动引理，便得到一个拟同构的 $\Ai$ 范畴。\znote{52}{扰动引理是这套证明的技术核心：若一个复形经收缩同伦等价于较小的复形，那么对它施加的“小扰动”可以转移到小复形上，并给出显式的级数公式。命题 3.8 的证明正是先用 $T^1$ 把 $\widetilde{\mathcal D}$ 的微分扰动为 $\mu^1_{\widetilde{\mathcal D}}$，再由扰动引理读出所需的 $\Ai$ 范畴。[来源：nLab "homological perturbation theory"，https:/\allowbreak /\allowbreak ncatlab.org/\allowbreak nlab/\allowbreak show/\allowbreak homological+\allowbreak perturbation+\allowbreak theory（2026-11-15 访问）]}显式计算表明，它恰好就是\stref{命题}{3.8}中的范畴 $\mathcal D$（注意，出于次数的原因，阶数大于三的复合都为零，故计算到此为止）。因此得到 $\Ai$ 函子 $\G:\mathcal D\to\widetilde{\mathcal D}\subset\Tw\A$，满足
\[
\G(Z_k)=\widetilde Z_k,\quad[\G^1(x_1)]=[c],\quad[\G^1(x_2)]=[i],\quad[\G^1(x_3)]=[p].
\]

现在设 $H(\A)$ 中有正合三角\eref{3.9}。令 $\C\subset\A$ 为由三角的三个对象组成的全子范畴；令 $\widetilde{\C}\subset\Tw\A$ 为包含这些对象在 $\A\hookrightarrow\Tw\A$ 下的像以及 $\Cone(c)$ 的全子范畴，其中 $c=c_1$。由\stref{引理}{3.31}与\stref{引理}{3.27}，显然的嵌入 $\C\hookrightarrow\widetilde{\C}$ 是拟等价，故由\stref{定理}{2.9}可以找到逆向拟等价 $\mathcal K:\widetilde{\C}\to\C$。这个逆在对象 $Y_0,Y_1$ 上为恒等，将 $\Cone(c)$ 映为 $Y_2$，并将 $[i],[p]$ 分别映为 $[c_2],[c_3]$。于是 $\F=\mathcal K\circ\G:\mathcal D\to\C\subset\A$ 就是具有所需性质的 $\Ai$ 函子。

还需去掉严格含幺的假设。设 $\A$ 仅为 c-含幺，并考虑其中的某个正合三角。先取最小模型，再用\stref{引理}{2.1}将其严格化，得到严格含幺范畴 $\widetilde{\A}$ 及拟等价 $\Phi:\A\to\widetilde{\A}$。由\stref{推论}{3.9}，所给三角在 $\Phi$ 下的像仍正合。对这个像应用前面的讨论，得到具有所需性质的 $\Ai$ 函子 $\widetilde{\F}:\mathcal D\to\widetilde{\A}$。再取\stref{定理}{2.9}给出的逆向拟等价 $\Psi$，令 $\F=\Psi\circ\widetilde{\F}$，并利用其与原对象及态射的上同调同构作识别，即得结论。
\end{proof}

\begin{lemma}\label{st:3.32}
$\Tw\A$ 由其全子范畴 $\A$ 生成。
\end{lemma}
\begin{proof}
取扭曲复形 $Y$ 及其子复形滤过 $F^iY$。由于 $\delta_Y$ 严格下三角，对每个 $i$，它都诱导映射 $d^i\in\hom^1_{\Tw\A}(F^iY/F^{i+1}Y,F^{i+1}Y)$。若将其视为从 $S^{-1}(F^iY/F^{i+1}Y)$ 到 $F^{i+1}Y$ 的零次态射，那么相应映射锥同构于 $F^iY$，故有正合三角
\SourceEquationDiagram{diagram-3-31}{0.31361505}{3.31}
\SourcePage{61}
由定义，$F^iY/F^{i+1}Y$ 的微分为零，因而同构于 $\A$ 中对象的若干平移的直和。利用这一点及\eref{3.31}，可见 $Y$ 能够由 $\A$ 中的对象反复施加锥与平移构造出来，这正是生成子范畴的定义所要求的。
\end{proof}

\begin{lemma}\label{st:3.33}
$\Ai$ 范畴 $\A$ 三角化，当且仅当嵌入 $\A\hookrightarrow\Tw\A$ 是拟等价。
\end{lemma}
\begin{proof}
设该嵌入为拟等价。$\A$ 中任意零次上闭链的抽象映射锥，都由 $\Tw\A$ 的某个对象拟表示，故由假设也由 $\A$ 的某个对象拟表示。任一方向的平移也具有同样性质；由定义，这就是说 $\A$ 是三角化的 $\Ai$ 范畴。

反之，设 $\A$ 三角化。由\stref{引理}{3.32}，$H^0(\Tw\A)$ 的任意对象，在同构意义下，都可以从 $H^0(\A)$ 的对象出发，反复取映射锥与任一方向的平移构造出来。不过，由于 $\A$ 三角化，所有这些运算都不会离开 $H^0(\Tw\A)$ 中的如下全子范畴：其对象同构于 $H^0(\A)$ 的某个对象。因此该嵌入是拟等价。
\end{proof}
由此结合\stref{命题}{3.29}与\stref{引理}{3.31}，可以推出\stref{命题}{3.14}。具体地，$H^0(\Tw\A)$ 是经典意义下的三角范畴；而 $H^0(\A)$ 是与之等价的全子范畴，具有相同的平移函子及相应的殊三角类，因此它本身也是三角范畴。\znote{53}{此处 $H^0(\A)$ 作为 $H^0(\Tw\A)$ 的全子范畴，与后者有相同的平移函子，殊三角类也相同（由\stref{引理}{3.31}），所以“本身也是三角范畴”并非形式推论：等价的全子范畴只要在三角运算下封闭，就继承三角结构。[来源：nLab "triangulated category"，https:/\allowbreak /\allowbreak ncatlab.org/\allowbreak nlab/\allowbreak show/\allowbreak triangulated+\allowbreak category（2026-11-15 访问）]}

\begin{lemma}\label{st:3.34}
设 $\F:\A\to\B$ 是上同调全忠实的 $\Ai$ 函子。假设 $\B$ 是三角化的 $\Ai$ 范畴，且 $\F$ 的像中的对象生成 $\B$。那么存在拟等价 $\widetilde{\F}:\Tw\A\to\B$，它在 $\A\subset\Tw\A$ 上的限制，在 $H^0(\operatorname{fun}(\A,\B))$ 中与 $\F$ 同构。
\end{lemma}
\begin{proof}
由假设，$\B\hookrightarrow\Tw\B$ 是拟等价。取\stref{定理}{2.9}给出的逆向拟等价 $\mathcal K:\Tw\B\to\B$，定义 $\widetilde{\F}=\mathcal K\circ\Tw\F$。它在 $\A$ 上的限制显然满足要求，而且至少是上同调全忠实的。由假设，$H^0(\B)$ 的任意对象，在同构意义下，都可以从形如 $\F(X)$ 的对象出发，反复取映射锥与任一方向的平移构造出来。但 $\widetilde{\F}$ 本身保持正合三角和平移，见\stref{推论}{3.9}，故所有这些运算都不会离开同构于 $\widetilde{\F}$ 的像中对象的那一类对象。因此 $\widetilde{\F}$ 是拟等价。
\end{proof}
这一结论结合\stref{引理}{3.32}，蕴含\stref{引理}{3.20}。

\SourcePage{62}
\subsection{\texorpdfstring{$\Ai$ 模的映射锥}{A∞ 模的映射锥}}\label{sec:3s}
现在考虑\stref{例题}{3.13}。设 $t:\M_0\to\M_1$ 是 $\Ai$ 模之间的零次同态。
\begin{definition*}
定义 $\C=\ACone(t)$ 为如下 $\Ai$ 模：
\begin{equation}\label{eq:3.32}
\begin{aligned}
\C(X)&=\M_0(X)[1]\oplus\M_1(X),\\
\mu^d_{\C}((b_0,b_1),a_{d-1},\ldots,a_1)
&=\bigl(\mu^d_{\M_0}(b_0,a_{d-1},\ldots,a_1),\\[-2pt]
&\qquad\mu^d_{\M_1}(b_1,a_{d-1},\ldots,a_1)+t^d(b_0,a_{d-1},\ldots,a_1)\bigr).
\end{aligned}
\end{equation}
\end{definition*}
这显然推广了\eref{3.6}；后者对应 $\M_k=\mathcal I(Y_k)$、$t=\mathcal I^1(c)$ 均来自 Yoneda 嵌入的情形。存在次数分别为零与一的典范模同态 $\iota:\M_1\to\C$、$\pi:\C\to\M_0$，其定义与刚才所说的特例完全相同。
\begin{lemma}\label{st:3.35}
$H^0(\Q)$ 中的下列三角正合：
\SourceDiagram{diagram-3-35}{0.21392285}
\end{lemma}
\begin{proof}
考虑全嵌入 $\Q\hookrightarrow\Tw\Q$。在后一个范畴中，有\eref{3.28}定义的典范锥 $\Cone(t)$。存在自然给出的态射
\begin{equation}\label{eq:3.33}
\C\longrightarrow\Cone(t),
\end{equation}
它对所有 $\M\in\Ob\Tw\Q$ 都诱导同构 $\hom_{\Tw\Q}(\M,\C)\to\hom_{\Tw\Q}(\M,\Cone(t))$（事实上，只需对 $\M\in\Ob\Q$ 检查这一点，因为一般情形可由一个直接的滤过论证推出）。这尤其意味着\eref{3.33}是同构。此外，很容易验证，它把 $\Cone(t)$ 与其组成部分之间的自然态射 $i,p$，同上面定义的 $\iota,\pi$ 联系起来。再用\stref{引理}{3.27}与\stref{引理}{3.31}，即得结论。
\end{proof}
这表明，$H^0(\Q)$ 中的任意态射都能补成一个正合三角。任一方向的平移模的存在性是显然的，因此 $\Q$ 三角化。\znote{54}{$\Q=\modcat(\A)$ 的三角化说明：$\Ai$ 模范畴可以充当 $\A$ 的三角化包络，这与导出范畴由链复形范畴对拟同构作局部化得到是同一类构造。区别在于这里用拟表示对象代替形式局部化，构造是具体的、可计算的。[来源：nLab "triangulated category"，https:/\allowbreak /\allowbreak ncatlab.org/\allowbreak nlab/\allowbreak show/\allowbreak triangulated+\allowbreak category（2026-11-15 访问）]}

为使讨论完整，我们回到\stref{笔记}{3.21}提出的要点。由于 $\Q$ 三角化，可以利用\stref{引理}{3.34}构造上同调全忠实函子 $\Tw\A\to\Q$。还有一种更具体的构造，参见\bcite{82}第 7 节：将 $\A$ 与 $\Tw\A$ 的 Yoneda 嵌入组合，得到 $\Ai$ 函子
\[
\widetilde{\mathcal I}:\Tw\A\xrightarrow{\mathcal I_{\Tw\A}}\modcat(\Tw\A)\xrightarrow{\mathcal I_{\A}^{*}}\modcat(\A)=\Q,
\]
其在 $\A\subset\Tw\A$ 上的限制显然等于 $\mathcal I_{\A}$。
\SourcePage{63}
\begin{lemma}\label{st:3.36}
$\widetilde{\mathcal I}$ 上同调全忠实。
\end{lemma}
\begin{proof}
考虑 $\Tw\A$ 中满足如下性质的对象 $Y$ 所成的子集：对所有 $X\in\Ob\A$，
\begin{equation}\label{eq:3.34}
H(\widetilde{\mathcal I}):\Hom_{H(\Tw\A)}(X,Y)\longrightarrow\Hom_{H(\Q)}(\widetilde{\mathcal I}(X),\widetilde{\mathcal I}(Y))
\end{equation}
都是同构。由于 $\widetilde{\mathcal I}$ 延伸了 $\A$ 的 Yoneda 嵌入，当 $Y$ 是 $\A$ 的对象时，这显然成立。\stref{推论}{3.9}说明，$\widetilde{\mathcal I}$ 将正合三角映为正合三角（事实上，可以直接验证：它严格地将锥映为锥，而不只是同构意义下如此）。利用\stref{推论}{3.15}中的相应长正合列，可以证明满足条件的 $Y$ 的类对映射锥封闭。再由平移相容性及\stref{引理}{3.32}，可知所有 $Y$ 都有这一性质。现在修改\eref{3.34}的要求，使双射性对所有 $X\in\Ob\Tw\A$ 都成立，再重复上述论证即可。
\end{proof}
